\section{연습문제}
\label{sec:field-extension-exercises}

문제는 A, B, C의 세 단계로 나뉜다. 별표 문제는 다음 절에 상세 풀이가 있다.

\subsection*{A. 정의와 계산}

\begin{exercise}
다음 체의 표수와 소체를 각각 구하라.
\[
  \Q,\quad \R,\quad \C,\quad \F_7,\quad \F_5(t),\quad \Q(t).
\]
\end{exercise}

\begin{exercise}
$\F_5$에서 프로베니우스 사상 $a\mapsto a^5$를 모든 원소에 대해 계산하라.
\end{exercise}

\begin{exercise}
원소가 $6$개, $12$개, $27$개인 체의 존재 가능성을 각각 판정하라. 여기서는 존재 가능한 원소 수의 형태만 판정하고, 실제 구성은 요구하지 않는다.
\end{exercise}

\begin{exercise}[*]
다음 확장의 차수와 자연스러운 기저를 구하라.
\begin{enumerate}[label=(\alph*)]
  \item $\Q(i)/\Q$,
  \item $\Q(\sqrt5)/\Q$,
  \item $\F_2(\alpha)/\F_2$, 단 $\alpha^2+\alpha+1=0$.
\end{enumerate}
\end{exercise}

\begin{exercise}
다음 최소다항식을 구하라.
\begin{enumerate}[label=(\alph*)]
  \item $1+\sqrt2$의 $\Q$ 위 최소다항식,
  \item $1/\sqrt2$의 $\Q$ 위 최소다항식,
  \item $i$의 $\R$ 위 최소다항식,
  \item $i$의 $\C$ 위 최소다항식.
\end{enumerate}
\end{exercise}

\begin{exercise}[*]
$\alpha=\sqrt[3]{2}$라 하자. $\Q(\alpha)$에서 다음을 표준형으로 계산하라.
\begin{enumerate}[label=(\alph*)]
  \item $\alpha^7$,
  \item $(1+\alpha)^{-1}$,
  \item $(1+\alpha+\alpha^2)^2$.
\end{enumerate}
\end{exercise}

\begin{exercise}
$K\subseteq E\subseteq L$이고 $[L:K]=18$이라 하자. 가능한 $[E:K]$의 값을 모두 적어라.
\end{exercise}

\begin{exercise}[*]
$\Q(\sqrt2)$에서 다음을 계산하라.
\begin{enumerate}[label=(\alph*)]
  \item $(3+2\sqrt2)^{-1}$,
  \item $(1+\sqrt2)^5$,
  \item $\sqrt2$의 $\Q(\sqrt2)$ 위 최소다항식.
\end{enumerate}
\end{exercise}

\begin{exercise}
$\F_2(\alpha)$에서 $\alpha^2+\alpha+1=0$이라 하자. 덧셈표와 곱셈표를 작성하고 영이 아닌 원소가 모두 역원을 가짐을 확인하라.
\end{exercise}

\begin{exercise}
$K(t)/K$에서 $1,t,t^2,\dots$가 $K$ 위에서 선형독립임을 보여 $[K(t):K]=\infty$임을 증명하라.
\end{exercise}

\subsection*{B. 핵심 증명}

\begin{exercise}
체의 표수가 $0$ 또는 소수임을 영인수를 이용하여 증명하라.
\end{exercise}

\begin{exercise}[*]
항등원을 보존하는 체 준동형 $\varphi:K\to L$가 존재하면 $\charac K=\charac L$임을 증명하라.
\end{exercise}

\begin{exercise}
표수가 $p>0$인 체에서 프로베니우스 준동형이 단사임을 두 가지 방법으로 증명하라: 체 준동형의 핵을 이용하는 방법과 $a^p=b^p\Rightarrow a=b$를 직접 보이는 방법.
\end{exercise}

\begin{exercise}[*]
유한체 $F$의 원소 수가 어떤 소수 $p$와 양의 정수 $n$에 대해 $p^n$임을 증명하라.
\end{exercise}

\begin{exercise}[*]
$K\subseteq L\subseteq M$이고 두 단계가 유한확장일 때, $L/K$의 기저와 $M/L$의 기저의 곱이 $M/K$의 기저임을 직접 증명하라.
\end{exercise}

\begin{exercise}
모든 유한확장이 대수적임을 거듭제곱들의 선형종속성을 이용하여 증명하라.
\end{exercise}

\begin{exercise}[*]
$\alpha$가 $K$ 위에서 대수적이고 최소다항식이 $m$이라 하자. $f(\alpha)=0$인 모든 $f\in K[X]$에 대해 $m\mid f$임을 나눗셈 알고리즘으로 증명하라.
\end{exercise}

\begin{exercise}[*]
$\alpha$가 $K$ 위에서 대수적이면 $K[\alpha]=K(\alpha)$임을 증명하라. 특히 영이 아닌 $f(\alpha)$의 역원이 다시 $K[\alpha]$에 속함을 베주 항등식으로 보여라.
\end{exercise}

\begin{exercise}
$\alpha_1,\dots,\alpha_r$가 모두 $K$ 위에서 대수적이면 $K(\alpha_1,\dots,\alpha_r)/K$가 유한확장임을 증명하라.
\end{exercise}

\begin{exercise}[*]
$K\subseteq L\subseteq M$에서 $L/K$와 $M/L$이 대수적이면 $M/K$도 대수적임을 증명하라.
\end{exercise}

\begin{exercise}
확장 $L/K$ 안에서 $K$ 위 대수적인 원소 전체가 $L$의 부분체를 이룸을 증명하라.
\end{exercise}

\begin{exercise}
$[L:K]=p$가 소수이면 $K$와 $L$ 사이에 진정한 중간체가 없음을 증명하라.
\end{exercise}

\subsection*{C. 구조와 응용}

\begin{exercise}[*]
$\Q(\sqrt2,\sqrt3)/\Q$의 차수가 $4$임을 증명하고 기저를 구하라.
\end{exercise}

\begin{exercise}[*]
$\alpha=\sqrt2+\sqrt3$의 $\Q$ 위 최소다항식을 구하라. 또한 $\Q(\alpha)=\Q(\sqrt2,\sqrt3)$임을 증명하라.
\end{exercise}

\begin{exercise}[*]
$\alpha=\sqrt[3]{2}$라 하자.
\begin{enumerate}[label=(\alph*)]
  \item $\alpha^2$의 $\Q$ 위 최소다항식을 구하라.
  \item $\Q(\alpha^2)=\Q(\alpha)$임을 보여라.
  \item $2+\alpha+\alpha^2$의 역원을 구하라.
\end{enumerate}
\end{exercise}

\begin{exercise}[*]
$f=X^3+X+1\in\F_2[X]$라 하고 $\alpha=X+(f)$라 하자.
\begin{enumerate}[label=(\alph*)]
  \item $f$가 기약임을 보여라.
  \item $\F_2(\alpha)$의 원소 수와 기저를 구하라.
  \item $\alpha^{-1}$과 $(\alpha+1)^{-1}$을 구하라.
  \item $\alpha^7=1$임을 확인하라.
\end{enumerate}
\end{exercise}

\begin{exercise}[*]
$K=\F_p(t)$에서 프로베니우스 준동형이 전사가 아님을 증명하라. 특히 $t$가 $K$ 안의 $p$제곱이 아님을 차수 합동으로 보여라.
\end{exercise}

\begin{exercise}
$[L:K]=2$이고 $\alpha\in L\setminus K$라 하자. $L=K(\alpha)$이고 $\alpha$의 최소다항식 차수가 $2$임을 증명하라.
\end{exercise}

\begin{exercise}
$[L:K]$가 홀수이면 $L/K$에는 $K$ 위 차수가 $2$인 중간체가 존재하지 않음을 증명하라.
\end{exercise}

\begin{exercise}[*]
$\alpha=\sqrt[3]{2}$이고 $\omega$가 $\omega^2+\omega+1=0$, $\omega\ne1$을 만족하는 복소수라 하자. 다음을 보여라.
\begin{enumerate}[label=(\alph*)]
  \item $[\Q(\alpha):\Q]=3$,
  \item $\omega\notin\Q(\alpha)$,
  \item $[\Q(\alpha,\omega):\Q]=6$,
  \item $1,\alpha,\alpha^2,\omega,\alpha\omega,\alpha^2\omega$가 $\Q$-기저이다.
\end{enumerate}
\end{exercise}
